% Version 2 worked-solutions exam, adapted from G10_T1_C1C2C3C4C5C6C7_exam_solution.tex.
% New contexts and data; assessment rules and mathematical models retained.
% Full-width worked solutions; monetary units are local to each problem.
\makeatletter
\def\input@path{{./}{../}{../../}{../../../}{../../../../}{preamble/}{../preamble/}{../../preamble/}{../../../preamble/}}
\makeatother
\input{preamble_exams.tex}
% Local solution-document spacing and readable labels in the inherited table.
\linespread{1.08}
\setlength{\parskip}{4pt}
% Retain the shared four-argument helper and its list handling, with a wider
% label column and explicit spaces for labels defined under ExplSyntaxOn.
\ExplSyntaxOn
\RenewDocumentCommand{\DecisionElementsTable}{mmmm}{
\decision_prepare_lists:nn {#1}{#2}
\begin{center}
\begin{tabularx}{\textwidth}{@{}p{2.6cm}>{\raggedright\arraybackslash}X@{}}
\toprule Element & Description \\ \midrule
Alternatives & \decision_print_slots:Nn \l_decision_alternatives_seq {A} \\
States~of~nature & \decision_print_slots:Nn \l_decision_states_seq {S} \\
Events & #3 \\ Consequences & #4 \\ \bottomrule
\end{tabularx}\end{center}
}
\ExplSyntaxOff
\title{Grade 10 -- Term 1\\C1--C7 Catch-up Exam Solutions -- Version 2}
\author{}
\date{2026--2027}
\begin{document}
\maketitle
\vspace{-1.5cm}
\setcounter{tocdepth}{1}
\setcounter{secnumdepth}{2}
\phantomsection\label{toc}
\tableofcontents
\clearpage
\begingroup\small
\section*{Learning objective}
Characterise economic decisions under uncertainty, organise payoff tables, calculate profits from revenues and costs, and justify choices using Maximax, Maximin, Minimax Regret, EMV/Bayes, and probability sensitivity analysis.

\section*{Criteria assessment}
This exam assesses C1--C7. Each question and task shows its points.
\textbf{Passing requires a minimum of 5 cumulative points per criterion,
including points earned in previous learning activities}.
Show all procedures; answers without the required procedure receive no credit.
Partial credit may be awarded for a correct method with arithmetic errors.
Missing names or late submissions incur a 1-point deduction per criterion.
\section*{Evaluated criteria}
\begin{enumerate}[label=\textbf{C\arabic*}]
\item Characterises decision-making problems under uncertainty by identifying decision alternatives, states of nature, consequences, and the economic or financial context of the decision.
\item Organises decision problems using payoff tables, placing alternatives, states of nature, and payoffs in a clear structure for comparison.
\item Calculates profits by combining revenues with fixed and variable costs across different states of nature, determining the corresponding payoffs for each alternative, and organising the results in payoff tables for comparison.
\item Applies the Maximin and Maximax criteria to make optimistic or conservative decisions without probabilities, selecting the alternative with the highest possible payoff or the best worst-case, interpreting risks or limitations.
\item Constructs opportunity loss tables and applies the Minimax Regret criterion to select the alternative that minimises the maximum opportunity loss.
\item Calculates Expected Monetary Value using stated or estimated probabilities and applies the Bayes decision criterion to select the alternative with the highest expected payoff.
\item Uses subjective probabilities and sensitivity analysis to examine how changes in probability assumptions affect the preferred decision.
\end{enumerate}

\begin{center}
	\small\setlength{\tabcolsep}{5pt}
	\begin{tabular}{|c|c|c|c|c||c|c|c|c|}
		\hline
		\textbf{\textnumero}&\textbf{C1}&\textbf{C2}&\textbf{C3}&\textbf{C4}&
		\textbf{\textnumero}&\textbf{C5}&\textbf{C6}&\textbf{C7}\\\hline
		1&3&6&--&--&6&2&--&--\\\hline
		2&3&--&6&--&7&2&--&--\\\hline
		3&--&--&--&2&8&2&--&--\\\hline
		4&--&--&--&2&9&--&2&--\\\hline
		5&--&--&--&2&10&--&2&--\\\hline
		&&&&&11&--&2&--\\\hline
		&&&&&12&--&--&3\\\hline
		&&&&&13&--&--&3\\\hline
		\textbf{Tot.}&\textbf{6}&\textbf{6}&\textbf{6}&\textbf{6}&
		\textbf{Tot.}&\textbf{6}&\textbf{6}&\textbf{6}\\\hline
	\end{tabular}
\end{center}

Each criterion has 6 available points; the exam has 42 available points overall.
\textbf{Notation.} $A_i$ and $S_j$ always follow the original alternative and state order within each problem. $P(A_i,S_j)$ denotes a monetary payoff; $P(S_j)$ denotes a probability. Navigation arrows return to the contents, the problem, or the summary.
\endgroup

\clearpage
\sectionScored{Community Science Sessions}{C1-\sml{3} C2-\sml{6}}{sec:community-science-sessions}
A community science centre must choose an astronomy session or a robotics session before family demand is known. Demand will be busy or quiet. Astronomy provides a net contribution of 14 USD per participant and a fixed sponsorship contribution of 90 USD per session. Robotics provides 17 USD per participant and a fixed sponsorship contribution of 65 USD per session. Busy demand adds a further net contribution of 5 USD per participant to either session; quiet demand adds 3 USD per participant. Astronomy attracts 28 participants with busy demand or 17 with quiet demand; robotics attracts 23 or 14, respectively. All contributions are additive, and no other adjustments are needed.
\QuestionsTasks[sec:community-science-sessions]{
{1}/{sol:community-science-sessions-c1}/{Characterize the key elements of the decision-making situation \sscore{C1-3}},
{2}/{sol:community-science-sessions-c2}/{and construct the payoff table. \sscore{C2-6}}
}
\subsection*{C1 - Characterization \score{3}}
\phantomsection\label{sol:community-science-sessions-c1}
\DecisionElementsTable{Astronomy session; Robotics session}
{Busy family demand; Quiet family demand}
{Family demand revealed after the session is chosen; the centre cannot control it.}
{Total net contribution in USD for each session--demand pair.}
The economic context is a community science centre choosing its session before family demand is known.
The decision controls the session, while demand controls attendance and the extra contribution rate.
\subsection*{C2 - Additive contributions and payoff table \score{6}}
\phantomsection\label{sol:community-science-sessions-c2}
Let $q_{ij}$ be participants, $v_i$ the session contribution per participant,
$d_j$ the demand contribution per participant, and $f_i$ the fixed sponsorship.
All monetary rates are in USD per participant; $f_i$ and payoffs are in USD.
\[
P(A_i,S_j)=q_{ij}(v_i+d_j)+f_i.
\]
The per-participant contributions are added before multiplying by participants.
The fixed sponsorship is then added once per session; it is not multiplied by participants.
\begin{align*}
P(A_1,S_1)&=28(14+5)+90=532+90=622,\\
P(A_1,S_2)&=17(14+3)+90=289+90=379,\\
P(A_2,S_1)&=23(17+5)+65=506+65=571,\\
P(A_2,S_2)&=14(17+3)+65=280+65=345.
\end{align*}
\textbf{Completed payoff table (USD).}
\begin{center}
\begin{tabular}{@{}lcc@{}}\toprule
Alternative & Busy ($S_1$) & Quiet ($S_2$)\\\midrule
Astronomy session ($A_1$) & 622 & 379\\
Robotics session ($A_2$) & 571 & 345\\
\bottomrule\end{tabular}
\end{center}
\textbf{Conclusion:} the astronomy session gives a higher contribution in both states: USD 622 under busy demand and USD 379 under quiet demand.
\par\vspace{4mm}
\secInicio[sec:community-science-sessions][sec:summarized-solutions]

\clearpage
\sectionScored{Kayak Rentals}{C1-\sml{3} C3-\sml{6}}{sec:kayak-rentals}
\begingroup\setlength{\parskip}{2pt}\linespread{1.02}\selectfont
A kayak rental business must choose single kayaks or tandem kayaks before seasonal demand is known. Demand will be busy or quiet. Single kayaks generate 84 rentals with busy demand or 56 with quiet demand; tandem kayaks generate 66 or 38, respectively. Base revenue per rental is 36 USD for the single kayaks and 49 USD for the tandem kayaks. Busy demand adds a surcharge of 7 USD per rental to either rental type; quiet demand adds 4 USD per rental. Variable cost per rental is 20 USD for the single kayaks and 29 USD for the tandem kayaks. Total fixed seasonal costs, for either rental type, are 540 USD with busy demand and 620 USD with quiet demand. These fixed costs are paid once per season, and there are no other revenues or costs.
\QuestionsTasks[sec:kayak-rentals]{
{1}/{sol:kayak-rentals-c1}/{Characterize the key elements of the decision-making situation \sscore{C1-3}},
{3}/{sol:kayak-rentals-c3}/{and construct the payoff table. \sscore{C3-6}}
}
\subsection*{C1 - Characterization \score{3}}
\phantomsection\label{sol:kayak-rentals-c1}
\DecisionElementsTable{Single kayaks; Tandem kayaks}
{Busy seasonal demand; Quiet seasonal demand}
{Seasonal demand revealed after the rental type is chosen; the operator cannot control it.}
{Seasonal profit in USD for each rental type--demand pair.}
The economic context is a kayak rental business choosing a rental type before knowing sales,
the demand surcharge, and the seasonal fixed cost.
\subsection*{C3 - Profit calculations and payoff table \score{6}}
\phantomsection\label{sol:kayak-rentals-c3}
Profit equals total revenue minus total variable costs minus fixed seasonal costs:
\begin{align*}
R_{ij}&=q_{ij}(b_i+s_j), & VC_{ij}&=q_{ij}c_i,\\
P(A_i,S_j)&=R_{ij}-VC_{ij}-F_j
=q_{ij}(b_i+s_j)-q_{ij}c_i-F_j.
\end{align*}
Here $q_{ij}$ is the number of rentals sold; $b_i$, $s_j$, and $c_i$ are base revenue, surcharge,
and variable cost in USD per rental; $F_j$ is the fixed seasonal cost in USD.
The surcharge is part of revenue for every rental sold. The fixed cost is paid
once in each alternative--state calculation, not once per rental.
\begin{align*}
P(A_1,S_1)&=84(36+7)-84(20)-540=3612-1680-540=1392,\\
P(A_1,S_2)&=56(36+4)-56(20)-620=2240-1120-620=500,\\
P(A_2,S_1)&=66(49+7)-66(29)-540=3696-1914-540=1242,\\
P(A_2,S_2)&=38(49+4)-38(29)-620=2014-1102-620=292.
\end{align*}
\textbf{Completed payoff table (profit in USD).}
\begin{center}
\begin{tabular}{@{}lcc@{}}\toprule
Alternative & Busy ($S_1$) & Quiet ($S_2$)\\\midrule
Single kayaks ($A_1$) & 1392 & 500\\
Tandem kayaks ($A_2$) & 1242 & 292\\
\bottomrule\end{tabular}
\end{center}
\textbf{Conclusion:} single kayaks earn more in both states: USD 1,392 under busy demand and USD 500 under quiet demand.
\par\vspace{4mm}
\secInicio[sec:kayak-rentals][sec:summarized-solutions]
\endgroup


\clearpage
\sectionScored{Bicycle Assembly}{C4-\sml{2}}{sec:bicycle-assembly}
A bicycle manufacturer must select an assembly plan before market conditions are known. Licensed production earns a fixed royalty regardless of the market state. The possible states, in column order, are Boom, Stable market, Recession. Profits are in thousand USD. No probabilities are available.
\begin{center}
\begin{tabular}{@{}lccc@{}}\toprule
Alternative & $S_1$ & $S_2$ & $S_3$\\\midrule
Large factory ($A_1$) & 132 & 53 & -27\\
Standard line ($A_2$) & 91 & 63 & 19\\
Compact line ($A_3$) & 61 & 47 & 33\\
Licensed production ($A_4$) & 26 & 26 & 26\\
\bottomrule\end{tabular}
\end{center}
\QuestionsTasks[sec:bicycle-assembly]{
{4}/{sol:bicycle-assembly-c4}/{What would the optimistic decision be? What would the conservative decision be? Justify your answer. \score{2}}
}
\subsection*{C4 - Maximax and Maximin \score{2}}
\phantomsection\label{sol:bicycle-assembly-c4}
Alternatives: $A_1$: Large factory; $A_2$: Standard line; $A_3$: Compact line; $A_4$: Licensed production.\par
States: $S_1$: Boom; $S_2$: Stable market; $S_3$: Recession.
All profits are in thousand USD. Find the largest and smallest payoff in every row:
\begin{center}
\begin{tabular}{@{}lcc@{}}\toprule
Alternative & Row maximum $M_i$ & Row minimum $m_i$\\\midrule
Large factory ($A_1$) & $\max\{132,53,-27\}=132$ & $\min\{132,53,-27\}=-27$\\
Standard line ($A_2$) & $\max\{91,63,19\}=91$ & $\min\{91,63,19\}=19$\\
Compact line ($A_3$) & $\max\{61,47,33\}=61$ & $\min\{61,47,33\}=33$\\
Licensed production ($A_4$) & $\max\{26,26,26\}=26$ & $\min\{26,26,26\}=26$\\
\bottomrule\end{tabular}
\end{center}
\[\max_i M_i=\max\{132,91,61,26\}=132,\qquad \max_i m_i=\max\{-27,19,33,26\}=33.\]
\textbf{Maximax:} choose \textbf{Large factory} ($A_1$), with best-case profit 132 thousand USD.\par
\textbf{Maximin:} choose \textbf{Compact line} ($A_3$), with worst-case profit 33 thousand USD.
Maximax focuses on the best possible profit and ignores the downside. Maximin protects the best worst-case profit and may forgo a large gain. Without probabilities, neither criterion estimates likelihoods or expected profit.
\par\vspace{4mm}
\secInicio[sec:bicycle-assembly][sec:summarized-solutions]

\clearpage
\sectionScored{Solar Installation}{C4-\sml{2}}{sec:solar-installation}
A solar installer must select an installation programme before customer adoption is known. The service agreement pays a fixed profit under every adoption state. The possible states, in column order, are Rapid adoption, Moderate adoption, Slow adoption. Profits are in thousand USD. No probabilities are available.
\begin{center}
\begin{tabular}{@{}lccc@{}}\toprule
Alternative & $S_1$ & $S_2$ & $S_3$\\\midrule
Regional rollout ($A_1$) & 168 & 51 & -39\\
Town rollout ($A_2$) & 108 & 69 & 17\\
Village rollout ($A_3$) & 73 & 52 & 31\\
Service agreement ($A_4$) & 35 & 35 & 35\\
\bottomrule\end{tabular}
\end{center}
\QuestionsTasks[sec:solar-installation]{
{4}/{sol:solar-installation-c4}/{What would the optimistic decision be? What would the conservative decision be? Justify your answer. \score{2}}
}
\subsection*{C4 - Maximax and Maximin \score{2}}
\phantomsection\label{sol:solar-installation-c4}
Alternatives: $A_1$: Regional rollout; $A_2$: Town rollout; $A_3$: Village rollout; $A_4$: Service agreement.\par
States: $S_1$: Rapid adoption; $S_2$: Moderate adoption; $S_3$: Slow adoption.
All profits are in thousand USD. Find the largest and smallest payoff in every row:
\begin{center}
\begin{tabular}{@{}lcc@{}}\toprule
Alternative & Row maximum $M_i$ & Row minimum $m_i$\\\midrule
Regional rollout ($A_1$) & $\max\{168,51,-39\}=168$ & $\min\{168,51,-39\}=-39$\\
Town rollout ($A_2$) & $\max\{108,69,17\}=108$ & $\min\{108,69,17\}=17$\\
Village rollout ($A_3$) & $\max\{73,52,31\}=73$ & $\min\{73,52,31\}=31$\\
Service agreement ($A_4$) & $\max\{35,35,35\}=35$ & $\min\{35,35,35\}=35$\\
\bottomrule\end{tabular}
\end{center}
\[\max_i M_i=\max\{168,108,73,35\}=168,\qquad \max_i m_i=\max\{-39,17,31,35\}=35.\]
\textbf{Maximax:} choose \textbf{Regional rollout} ($A_1$), with best-case profit 168 thousand USD.\par
\textbf{Maximin:} choose \textbf{Service agreement} ($A_4$), with worst-case profit 35 thousand USD.
Maximax focuses on the best possible profit and ignores the downside. Maximin protects the best worst-case profit and may forgo a large gain. Without probabilities, neither criterion estimates likelihoods or expected profit.
\par\vspace{4mm}
\secInicio[sec:solar-installation][sec:summarized-solutions]

\clearpage
\sectionScored{Equipment Leasing}{C4-\sml{2}}{sec:equipment-leasing}
A construction-equipment business must choose a fleet plan before next season's rental utilisation is known. The possible states, in column order, are Peak utilisation, Busy utilisation, Routine utilisation, Quiet utilisation. Profits are in thousand USD. No probabilities are available.
\begin{center}
\begin{tabular}{@{}lcccc@{}}\toprule
Alternative & $S_1$ & $S_2$ & $S_3$ & $S_4$\\\midrule
Buy large fleet ($A_1$) & 156 & 81 & 21 & -46\\
Buy medium fleet ($A_2$) & 113 & 85 & 41 & 11\\
Lease small fleet ($A_3$) & 77 & 63 & 49 & 25\\
Pool equipment ($A_4$) & 54 & 51 & 39 & 32\\
\bottomrule\end{tabular}
\end{center}
\QuestionsTasks[sec:equipment-leasing]{
{4}/{sol:equipment-leasing-c4}/{What would the optimistic decision be? What would the conservative decision be? Justify your answer. \score{2}}
}
\subsection*{C4 - Maximax and Maximin \score{2}}
\phantomsection\label{sol:equipment-leasing-c4}
Alternatives: $A_1$: Buy large fleet; $A_2$: Buy medium fleet; $A_3$: Lease small fleet; $A_4$: Pool equipment.\par
States: $S_1$: Peak utilisation; $S_2$: Busy utilisation; $S_3$: Routine utilisation; $S_4$: Quiet utilisation.
All profits are in thousand USD. Find the largest and smallest payoff in every row:
\begin{center}
\begin{tabular}{@{}lcc@{}}\toprule
Alternative & Row maximum $M_i$ & Row minimum $m_i$\\\midrule
Buy large fleet ($A_1$) & $\max\{156,81,21,-46\}=156$ & $\min\{156,81,21,-46\}=-46$\\
Buy medium fleet ($A_2$) & $\max\{113,85,41,11\}=113$ & $\min\{113,85,41,11\}=11$\\
Lease small fleet ($A_3$) & $\max\{77,63,49,25\}=77$ & $\min\{77,63,49,25\}=25$\\
Pool equipment ($A_4$) & $\max\{54,51,39,32\}=54$ & $\min\{54,51,39,32\}=32$\\
\bottomrule\end{tabular}
\end{center}
\[\max_i M_i=\max\{156,113,77,54\}=156,\qquad \max_i m_i=\max\{-46,11,25,32\}=32.\]
\textbf{Maximax:} choose \textbf{Buy large fleet} ($A_1$), with best-case profit 156 thousand USD.\par
\textbf{Maximin:} choose \textbf{Pool equipment} ($A_4$), with worst-case profit 32 thousand USD.
Maximax focuses on the best possible profit and ignores the downside. Maximin protects the best worst-case profit and may forgo a large gain. Without probabilities, neither criterion estimates likelihoods or expected profit.
\par\vspace{4mm}
\secInicio[sec:equipment-leasing][sec:summarized-solutions]

\clearpage
\sectionScored{Festival Meals}{C5-\sml{2}}{sec:festival-meals}
A festival caterer must select a meal-preparation plan before attendance is known. The possible states, in column order, are Crowded, Average attendance, Sparse attendance. Profits are in thousand USD. No probabilities are available.
\begin{center}
\begin{tabular}{@{}lccc@{}}\toprule
Alternative & $S_1$ & $S_2$ & $S_3$\\\midrule
Full menu ($A_1$) & 87 & 37 & -16\\
Balanced menu ($A_2$) & 65 & 54 & 21\\
Limited menu ($A_3$) & 43 & 41 & 34\\
\bottomrule\end{tabular}
\end{center}
\QuestionsTasks[sec:festival-meals]{
{5}/{sol:festival-meals-c5}/{What decision would be selected using the Minimax Regret criterion? Justify your answer. \score{2}}
}
\subsection*{C5 - Minimax Regret \score{2}}
\phantomsection\label{sol:festival-meals-c5}
Alternatives: $A_1$: Full menu; $A_2$: Balanced menu; $A_3$: Limited menu.\par
States: $S_1$: Crowded; $S_2$: Average attendance; $S_3$: Sparse attendance.
\[\boldsymbol{b}=\bigl(87,54,34\bigr),\qquad b_j=\max_k P(A_k,S_j).\]
Each column is compared with its highest profit. Regret is the profit missed by not choosing that state's best alternative:
\[r(A_i,S_j)=\max_k P(A_k,S_j)-P(A_i,S_j).\]
All payoffs and regrets in this calculation are in thousand USD.
\begin{center}
\begin{tabular}{@{}lcccc@{}}\toprule
Alternative & $S_1$ & $S_2$ & $S_3$ & Maximum regret\\\midrule
$A_1$ & $\begin{gathered}87-87\\=0\end{gathered}$ & $\begin{gathered}54-37\\=17\end{gathered}$ & $\begin{gathered}34-(-16)\\=50\end{gathered}$ & 50\\
$A_2$ & $\begin{gathered}87-65\\=22\end{gathered}$ & $\begin{gathered}54-54\\=0\end{gathered}$ & $\begin{gathered}34-21\\=13\end{gathered}$ & 22\\
$A_3$ & $\begin{gathered}87-43\\=44\end{gathered}$ & $\begin{gathered}54-41\\=13\end{gathered}$ & $\begin{gathered}34-34\\=0\end{gathered}$ & 44\\
\bottomrule\end{tabular}
\end{center}
\[\min_i\max_j r(A_i,S_j)=\min\{50,22,44\}=22.\]
Choose \textbf{Balanced menu} ($A_2$): its maximum regret is \textbf{22 thousand USD}. This is the smallest worst possible opportunity loss; it does not mean that this loss occurs in every state or that this alternative always earns the most profit.
\par\vspace{4mm}
\secInicio[sec:festival-meals][sec:summarized-solutions]

\clearpage
\sectionScored{Laundry Capacity}{C5-\sml{2}}{sec:laundry-capacity}
A commercial laundry must select its capacity plan before hotel orders are known. A contracted service provides a fixed profit under every workload state. The possible states, in column order, are Heavy workload, Routine workload, Light workload. Profits are in thousand USD. No probabilities are available.
\begin{center}
\begin{tabular}{@{}lccc@{}}\toprule
Alternative & $S_1$ & $S_2$ & $S_3$\\\midrule
Build plant ($A_1$) & 126 & 47 & -32\\
Upgrade machines ($A_2$) & 94 & 73 & 19\\
Lease machines ($A_3$) & 69 & 61 & 38\\
Contract service ($A_4$) & 41 & 41 & 41\\
\bottomrule\end{tabular}
\end{center}
\QuestionsTasks[sec:laundry-capacity]{
{5}/{sol:laundry-capacity-c5}/{What decision would be selected using the Minimax Regret criterion? Justify your answer. \score{2}}
}
\subsection*{C5 - Minimax Regret \score{2}}
\phantomsection\label{sol:laundry-capacity-c5}
Alternatives: $A_1$: Build plant; $A_2$: Upgrade machines; $A_3$: Lease machines; $A_4$: Contract service.\par
States: $S_1$: Heavy workload; $S_2$: Routine workload; $S_3$: Light workload.
\[\boldsymbol{b}=\bigl(126,73,41\bigr),\qquad b_j=\max_k P(A_k,S_j).\]
Each column is compared with its highest profit. Regret is the profit missed by not choosing that state's best alternative:
\[r(A_i,S_j)=\max_k P(A_k,S_j)-P(A_i,S_j).\]
All payoffs and regrets in this calculation are in thousand USD.
\begin{center}
\begin{tabular}{@{}lcccc@{}}\toprule
Alternative & $S_1$ & $S_2$ & $S_3$ & Maximum regret\\\midrule
$A_1$ & $\begin{gathered}126-126\\=0\end{gathered}$ & $\begin{gathered}73-47\\=26\end{gathered}$ & $\begin{gathered}41-(-32)\\=73\end{gathered}$ & 73\\
$A_2$ & $\begin{gathered}126-94\\=32\end{gathered}$ & $\begin{gathered}73-73\\=0\end{gathered}$ & $\begin{gathered}41-19\\=22\end{gathered}$ & 32\\
$A_3$ & $\begin{gathered}126-69\\=57\end{gathered}$ & $\begin{gathered}73-61\\=12\end{gathered}$ & $\begin{gathered}41-38\\=3\end{gathered}$ & 57\\
$A_4$ & $\begin{gathered}126-41\\=85\end{gathered}$ & $\begin{gathered}73-41\\=32\end{gathered}$ & $\begin{gathered}41-41\\=0\end{gathered}$ & 85\\
\bottomrule\end{tabular}
\end{center}
\[\min_i\max_j r(A_i,S_j)=\min\{73,32,57,85\}=32.\]
Choose \textbf{Upgrade machines} ($A_2$): its maximum regret is \textbf{32 thousand USD}. This is the smallest worst possible opportunity loss; it does not mean that this loss occurs in every state or that this alternative always earns the most profit.
\par\vspace{4mm}
\secInicio[sec:laundry-capacity][sec:summarized-solutions]

\clearpage
\sectionScored{Juice Production}{C5-\sml{2}}{sec:juice-production}
A juice producer must choose its daily production plan before sales are known. The possible states, in column order, are Exceptional sales, Brisk sales, Ordinary sales, Poor sales. Profits are in USD. No probabilities are available.
\begin{center}
\begin{tabular}{@{}lcccc@{}}\toprule
Alternative & $S_1$ & $S_2$ & $S_3$ & $S_4$\\\midrule
Large batch ($A_1$) & 1080 & 610 & 230 & -270\\
Staged batches ($A_2$) & 810 & 730 & 510 & 150\\
Small batch ($A_3$) & 530 & 550 & 480 & 390\\
\bottomrule\end{tabular}
\end{center}
\QuestionsTasks[sec:juice-production]{
{5}/{sol:juice-production-c5}/{What decision would be selected using the Minimax Regret criterion? Justify your answer. \score{2}}
}
\subsection*{C5 - Minimax Regret \score{2}}
\phantomsection\label{sol:juice-production-c5}
Alternatives: $A_1$: Large batch; $A_2$: Staged batches; $A_3$: Small batch.\par
States: $S_1$: Exceptional sales; $S_2$: Brisk sales; $S_3$: Ordinary sales; $S_4$: Poor sales.
\[\boldsymbol{b}=\bigl(1080,730,510,390\bigr),\qquad b_j=\max_k P(A_k,S_j).\]
Each column is compared with its highest profit. Regret is the profit missed by not choosing that state's best alternative:
\[r(A_i,S_j)=\max_k P(A_k,S_j)-P(A_i,S_j).\]
All payoffs and regrets in this calculation are in USD.
\begin{center}
\begin{tabular}{@{}lccccc@{}}\toprule
Alternative & $S_1$ & $S_2$ & $S_3$ & $S_4$ & Maximum regret\\\midrule
$A_1$ & $\begin{gathered}1080-1080\\=0\end{gathered}$ & $\begin{gathered}730-610\\=120\end{gathered}$ & $\begin{gathered}510-230\\=280\end{gathered}$ & $\begin{gathered}390-(-270)\\=660\end{gathered}$ & 660\\
$A_2$ & $\begin{gathered}1080-810\\=270\end{gathered}$ & $\begin{gathered}730-730\\=0\end{gathered}$ & $\begin{gathered}510-510\\=0\end{gathered}$ & $\begin{gathered}390-150\\=240\end{gathered}$ & 270\\
$A_3$ & $\begin{gathered}1080-530\\=550\end{gathered}$ & $\begin{gathered}730-550\\=180\end{gathered}$ & $\begin{gathered}510-480\\=30\end{gathered}$ & $\begin{gathered}390-390\\=0\end{gathered}$ & 550\\
\bottomrule\end{tabular}
\end{center}
\[\min_i\max_j r(A_i,S_j)=\min\{660,270,550\}=270.\]
Choose \textbf{Staged batches} ($A_2$): its maximum regret is \textbf{270 USD}. This is the smallest worst possible opportunity loss; it does not mean that this loss occurs in every state or that this alternative always earns the most profit.
\par\vspace{4mm}
\secInicio[sec:juice-production][sec:summarized-solutions]

\clearpage
\sectionScored{Coffee Roasting}{C6-\sml{2}}{sec:coffee-roasting}
A coffee roaster must choose a roasting plan before export conditions are known. The possible states, in column order, are Export boom, Regular exports, Export slump. Profits are in thousand USD.
\begin{center}
\begin{tabular}{@{}lccc@{}}\toprule
Alternative & $S_1$ & $S_2$ & $S_3$\\\midrule
$P(S_j)$ & 0.35 & 0.40 & 0.25\\
Specialty roast ($A_1$) & 98 & 49 & -17\\
Blended roast ($A_2$) & 74 & 63 & 27\\
Everyday roast ($A_3$) & 51 & 53 & 42\\
\bottomrule\end{tabular}
\end{center}
\QuestionsTasks[sec:coffee-roasting]{
{6}/{sol:coffee-roasting-c6}/{What decision would be selected using the EMV or Bayes criterion? Justify your answer. \score{2}}
}
\subsection*{C6 - EMV / Bayes \score{2}}
\phantomsection\label{sol:coffee-roasting-c6}
Alternatives: $A_1$: Specialty roast; $A_2$: Blended roast; $A_3$: Everyday roast.\par
States: $S_1$: Export boom; $S_2$: Regular exports; $S_3$: Export slump.
\[\boldsymbol{q}=(0.35,0.40,0.25),\qquad 0.35+0.40+0.25=1.\]
All probabilities lie in $[0,1]$, so they are nonnegative and valid.
\[EMV(A_i)=\sum_{j=1}^{n}P(S_j)P(A_i,S_j).\]
Multiply each state payoff by its probability, including any negative payoff, and add the weighted amounts. All EMVs are in thousand USD.
\begin{align*}
EMV(A_1)&=0.35(98)+0.40(49)+0.25(-17)\\&=34.30+19.60-4.25=49.65,\\
EMV(A_2)&=0.35(74)+0.40(63)+0.25(27)\\&=25.90+25.20+6.75=57.85,\\
EMV(A_3)&=0.35(51)+0.40(53)+0.25(42)\\&=17.85+21.20+10.50=49.55.
\end{align*}
\[\max\{49.65,57.85,49.55\}=57.85.\]
Choose \textbf{Blended roast} ($A_2$), with expected profit \textbf{57.85 thousand USD} (USD 57,850). EMV is a probability-weighted expected payoff, not a guaranteed outcome. The realised profit will be the payoff for the state that actually occurs.
\par\vspace{4mm}
\secInicio[sec:coffee-roasting][sec:summarized-solutions]

\clearpage
\sectionScored{Farm Visitor Facilities}{C6-\sml{2}}{sec:farm-visitor-facilities}
A farm must select visitor facilities before seasonal attendance is known. Leasing to an operator earns a fixed profit under all attendance states. The possible states, in column order, are Busy season, Typical season, Quiet season. Profits are in thousand USD.
\begin{center}
\begin{tabular}{@{}lccc@{}}\toprule
Alternative & $S_1$ & $S_2$ & $S_3$\\\midrule
$P(S_j)$ & 0.25 & 0.45 & 0.30\\
Visitor centre ($A_1$) & 117 & 43 & -29\\
Covered terrace ($A_2$) & 89 & 64 & 23\\
Picnic shelter ($A_3$) & 63 & 57 & 40\\
Lease to operator ($A_4$) & 32 & 32 & 32\\
\bottomrule\end{tabular}
\end{center}
\QuestionsTasks[sec:farm-visitor-facilities]{
{6}/{sol:farm-visitor-facilities-c6}/{What decision would be selected using the EMV or Bayes criterion? Justify your answer. \score{2}}
}
\subsection*{C6 - EMV / Bayes \score{2}}
\phantomsection\label{sol:farm-visitor-facilities-c6}
Alternatives: $A_1$: Visitor centre; $A_2$: Covered terrace; $A_3$: Picnic shelter; $A_4$: Lease to operator.\par
States: $S_1$: Busy season; $S_2$: Typical season; $S_3$: Quiet season.
\[\boldsymbol{q}=(0.25,0.45,0.30),\qquad 0.25+0.45+0.30=1.\]
All probabilities lie in $[0,1]$, so they are nonnegative and valid.
\[EMV(A_i)=\sum_{j=1}^{n}P(S_j)P(A_i,S_j).\]
Multiply each state payoff by its probability, including any negative payoff, and add the weighted amounts. All EMVs are in thousand USD.
\begin{align*}
EMV(A_1)&=0.25(117)+0.45(43)+0.30(-29)\\&=29.25+19.35-8.70=39.90,\\
EMV(A_2)&=0.25(89)+0.45(64)+0.30(23)\\&=22.25+28.80+6.90=57.95,\\
EMV(A_3)&=0.25(63)+0.45(57)+0.30(40)\\&=15.75+25.65+12=53.40,\\
EMV(A_4)&=0.25(32)+0.45(32)+0.30(32)\\&=8+14.40+9.60=32.
\end{align*}
\[\max\{39.90,57.95,53.40,32\}=57.95.\]
Choose \textbf{Covered terrace} ($A_2$), with expected profit \textbf{57.95 thousand USD} (USD 57,950). EMV is a probability-weighted expected payoff, not a guaranteed outcome. The realised profit will be the payoff for the state that actually occurs.
\par\vspace{4mm}
\secInicio[sec:farm-visitor-facilities][sec:summarized-solutions]

\clearpage
\sectionScored{Translation Capacity}{C6-\sml{2}}{sec:translation-capacity}
A translation agency must choose a capacity plan before client orders are known. The possible states, in column order, are Order rush, Busy orders, Routine orders, Few orders. Profits are in thousand USD.
\begin{center}
\begin{tabular}{@{}lcccc@{}}\toprule
Alternative & $S_1$ & $S_2$ & $S_3$ & $S_4$\\\midrule
$P(S_j)$ & 0.15 & 0.35 & 0.25 & 0.25\\
Permanent staff ($A_1$) & 137 & 81 & 29 & -43\\
Staff plus freelancers ($A_2$) & 99 & 86 & 59 & 17\\
Core staff ($A_3$) & 65 & 67 & 61 & 48\\
\bottomrule\end{tabular}
\end{center}
\QuestionsTasks[sec:translation-capacity]{
{6}/{sol:translation-capacity-c6}/{What decision would be selected using the EMV or Bayes criterion? Justify your answer. \score{2}}
}
\subsection*{C6 - EMV / Bayes \score{2}}
\phantomsection\label{sol:translation-capacity-c6}
Alternatives: $A_1$: Permanent staff; $A_2$: Staff plus freelancers; $A_3$: Core staff.\par
States: $S_1$: Order rush; $S_2$: Busy orders; $S_3$: Routine orders; $S_4$: Few orders.
\[\boldsymbol{q}=(0.15,0.35,0.25,0.25),\qquad 0.15+0.35+0.25+0.25=1.\]
All probabilities lie in $[0,1]$, so they are nonnegative and valid.
\[EMV(A_i)=\sum_{j=1}^{n}P(S_j)P(A_i,S_j).\]
Multiply each state payoff by its probability, including any negative payoff, and add the weighted amounts. All EMVs are in thousand USD.
\begin{align*}
EMV(A_1)&=0.15(137)+0.35(81)+0.25(29)+0.25(-43)\\&=20.55+28.35+7.25-10.75=45.40,\\
EMV(A_2)&=0.15(99)+0.35(86)+0.25(59)+0.25(17)\\&=14.85+30.10+14.75+4.25=63.95,\\
EMV(A_3)&=0.15(65)+0.35(67)+0.25(61)+0.25(48)\\&=9.75+23.45+15.25+12=60.45.
\end{align*}
\[\max\{45.40,63.95,60.45\}=63.95.\]
Choose \textbf{Staff plus freelancers} ($A_2$), with expected profit \textbf{63.95 thousand USD} (USD 63,950). EMV is a probability-weighted expected payoff, not a guaranteed outcome. The realised profit will be the payoff for the state that actually occurs.
\par\vspace{4mm}
\secInicio[sec:translation-capacity][sec:summarized-solutions]

\clearpage
\sectionScored{Mobile Game Release}{C7-\sml{3}}{sec:mobile-game-release}
A game developer must choose a wide release or a limited release before player uptake is known. The possible states, in column order, are Viral uptake, Moderate uptake, Low uptake. Profits are in thousand USD.
\begin{center}
\begin{tabular}{@{}lccc@{}}\toprule
Alternative & $S_1$ & $S_2$ & $S_3$\\\midrule
Wide release ($A_1$) & 108 & 44 & -10\\
Limited release ($A_2$) & 58 & 64 & 30\\
\bottomrule\end{tabular}
\end{center}
Suppose $P(S_1)=p$ and $P(S_2)=p/2$; determine $P(S_3)$ from the sum-to-one condition.
\QuestionsTasks[sec:mobile-game-release]{
{7}/{sol:mobile-game-release-c7}/{State the preferred alternative for every probability value ($0\le p\le1$). \score{3}}
}
\subsection*{C7 - Probability sensitivity analysis \score{3}}
\phantomsection\label{sol:mobile-game-release-c7}
Alternatives: $A_1$: Wide release; $A_2$: Limited release.\par
States: $S_1$: Viral uptake; $S_2$: Moderate uptake; $S_3$: Low uptake.
The mutually exclusive and exhaustive state probabilities sum to one:
\[P(S_3)=1-p-p/2=1-3p/2,\qquad\boldsymbol{q}=\left(p,\frac p2,1-\frac{3p}{2}\right).\]
Every probability must lie in $[0,1]$:
\[0\le p\le1,\quad 0\le p/2\le1\iff0\le p\le2,\quad 0\le1-3p/2\le1\iff0\le p\le\frac23.\]
Their intersection is $\boxed{0\le p\le \frac{2}{3}}$.
The remaining part of $[0,1]$, $\frac{2}{3}<p\le1$, is infeasible because $P(S_3)<0$. No decision is valid there.
Expected profits are in thousand USD:
\begin{align*}
E_1&=EMV(A_1)=108p+44(p/2)-10(1-3p/2)\\&=108p+22p-10+15p=-10+145p,\\
E_2&=EMV(A_2)=58p+64(p/2)+30(1-3p/2)\\&=58p+32p+30-45p=30+45p.
\end{align*}
Equating the two linear EMVs gives
\[E_1=E_2\iff 100p=40\iff\boxed{p=\frac{2}{5}}.\]
This is the only threshold, since the slopes differ, and $0<\frac{2}{5}<\frac{2}{3}$. At the tie, $E_1=E_2=48$ thousand USD.
The difference $E_1-E_2=-40+100p$ is negative below the threshold and positive above it.
At $p=0$, $(E_1,E_2)=(-10,30)$; at $p=\frac{2}{3}$, $(E_1,E_2)=(\frac{260}{3},60)$.
Both feasible endpoints agree with the rule:
\[\boxed{\begin{cases}
\text{Limited release }(A_2),&0\le p<\frac{2}{5},\\
\text{Both alternatives tied},&p=\frac{2}{5},\\
\text{Wide release }(A_1),&\frac{2}{5}<p\le\frac{2}{3}.
\end{cases}}\]
Either alternative is optimal at the tie. Expected profit is not guaranteed profit.
\par\vspace{4mm}
\secInicio[sec:mobile-game-release][sec:summarized-solutions]

\clearpage
\sectionScored{Timber Supply Agreement}{C7-\sml{3}}{sec:timber-supply-agreement}
A timber merchant must choose spot sales or a reserved-supply agreement before next season's timber availability is known. The possible states, in column order, are Scarce timber, Balanced supply, Abundant timber. Profits are in thousand USD.
\begin{center}
\begin{tabular}{@{}lccc@{}}\toprule
Alternative & $S_1$ & $S_2$ & $S_3$\\\midrule
Spot sales ($A_1$) & 30 & 66 & 94\\
Reserved supply ($A_2$) & 85 & 46 & 40\\
\bottomrule\end{tabular}
\end{center}
Suppose $P(S_1)=2p$ and $P(S_2)=p$; determine $P(S_3)$ from the sum-to-one condition.
\QuestionsTasks[sec:timber-supply-agreement]{
{7}/{sol:timber-supply-agreement-c7}/{State the preferred alternative for every probability value ($0\le p\le1$). \score{3}}
}
\subsection*{C7 - Probability sensitivity analysis \score{3}}
\phantomsection\label{sol:timber-supply-agreement-c7}
Alternatives: $A_1$: Spot sales; $A_2$: Reserved supply.\par
States: $S_1$: Scarce timber; $S_2$: Balanced supply; $S_3$: Abundant timber.
The mutually exclusive and exhaustive state probabilities sum to one:
\[P(S_3)=1-2p-p=1-3p,\qquad\boldsymbol{q}=(2p,p,1-3p).\]
Every probability must lie in $[0,1]$:
\[0\le2p\le1\iff0\le p\le\frac12,\quad 0\le p\le1,\quad 0\le1-3p\le1\iff0\le p\le\frac13.\]
Their intersection is $\boxed{0\le p\le \frac{1}{3}}$.
The remaining part of $[0,1]$, $\frac{1}{3}<p\le1$, is infeasible because $P(S_3)<0$. No decision is valid there.
For $p>\frac12$, $P(S_1)>1$ as well.
Expected profits are in thousand USD:
\begin{align*}
E_1&=EMV(A_1)=30(2p)+66p+94(1-3p)\\&=60p+66p+94-282p=94-156p,\\
E_2&=EMV(A_2)=85(2p)+46p+40(1-3p)\\&=170p+46p+40-120p=40+96p.
\end{align*}
Equating the two linear EMVs gives
\[E_1=E_2\iff -252p=-54\iff\boxed{p=\frac{3}{14}}.\]
This is the only threshold, since the slopes differ, and $0<\frac{3}{14}<\frac{1}{3}$. At the tie, $E_1=E_2=\frac{424}{7}$ thousand USD.
The difference $E_1-E_2=54-252p$ is positive below the threshold and negative above it.
At $p=0$, $(E_1,E_2)=(94,40)$; at $p=\frac{1}{3}$, $(E_1,E_2)=(42,72)$.
Both feasible endpoints agree with the rule:
\[\boxed{\begin{cases}
\text{Spot sales }(A_1),&0\le p<\frac{3}{14},\\
\text{Both alternatives tied},&p=\frac{3}{14},\\
\text{Reserved supply }(A_2),&\frac{3}{14}<p\le\frac{1}{3}.
\end{cases}}\]
Either alternative is optimal at the tie. Expected profit is not guaranteed profit.
\par\vspace{4mm}
\secInicio[sec:timber-supply-agreement][sec:summarized-solutions]

\clearpage
\sectionScored{Summarized Solutions}{C1--C7}{sec:summarized-solutions}
{\small\setlength{\parskip}{2pt}
The numbered summaries follow the exam order. $A_i$ and $S_j$ use the original
row and column order. $E_i=EMV(A_i)$; $b_j$ is the best payoff in state $S_j$;
$r_{ij}=b_j-P(A_i,S_j)$. Units are stated for each problem.
\subsection{Community Science Sessions \hfill C1--3 C2--6}
\label{summary:community-science-sessions}
Alternatives: $A_1$: Astronomy session; $A_2$: Robotics session.\par
States: $S_1$: Busy demand; $S_2$: Quiet demand.
Uncertain event: actual demand after the choice. Context: community science centre selecting a session; consequences are net contributions in USD.
\[P_{ij}=q_{ij}(v_i+d_j)+f_i.\]
The fixed sponsorship is added once; the two per-participant rates multiply attendance.
\begin{center}
\begin{tabular}{@{}lcc@{}}\toprule
Payoff (USD) & $S_1$: busy & $S_2$: quiet\\\midrule
Astronomy ($A_1$) & $28(14+5)+90=622$ & $17(14+3)+90=379$\\
Robotics ($A_2$) & $23(17+5)+65=571$ & $14(17+3)+65=345$\\
\bottomrule\end{tabular}
\end{center}
Astronomy gives the higher contribution in both demand states.
\par\vspace{4mm}
\secInicio[sec:community-science-sessions][sec:summarized-solutions]
\subsection{Kayak Rentals \hfill C1--3 C3--6}
\label{summary:kayak-rentals}
Alternatives: $A_1$: Single kayaks; $A_2$: Tandem kayaks.\par
States: $S_1$: Busy demand; $S_2$: Quiet demand.
Uncertain event: actual demand after the choice. Context: kayak rental business selecting a rental type; consequences are seasonal profits in USD.
\[P_{ij}=q_{ij}(b_i+s_j)-q_{ij}c_i-F_j.\]
Revenue includes the surcharge; the seasonal fixed cost is subtracted once.
\begin{center}
\begin{tabular}{@{}lcc@{}}\toprule
Profit (USD) & $S_1$: busy & $S_2$: quiet\\\midrule
Single ($A_1$) & $84(36+7)-84(20)-540=1392$ & $56(36+4)-56(20)-620=500$\\
Tandem ($A_2$) & $66(49+7)-66(29)-540=1242$ & $38(49+4)-38(29)-620=292$\\
\bottomrule\end{tabular}
\end{center}
Single kayaks earn more in both demand states.
\par\vspace{4mm}
\secInicio[sec:kayak-rentals][sec:summarized-solutions]
\clearpage
\clearpage
\subsection{Bicycle Assembly \hfill C4--2}
\label{summary:bicycle-assembly}
Alternatives: $A_1$: Large factory; $A_2$: Standard line; $A_3$: Compact line; $A_4$: Licensed production.\par
States: $S_1$: Boom; $S_2$: Stable market; $S_3$: Recession.
\begin{center}
\begin{tabular}{@{}lcc@{}}\toprule
Alternative & Row maximum $M_i$ & Row minimum $m_i$\\\midrule
Large factory ($A_1$) & 132 & -27\\
Standard line ($A_2$) & 91 & 19\\
Compact line ($A_3$) & 61 & 33\\
Licensed production ($A_4$) & 26 & 26\\
\bottomrule\end{tabular}
\end{center}
\[\max_i M_i=\max\{132,91,61,26\}=132,\qquad \max_i m_i=\max\{-27,19,33,26\}=33.\]
\textbf{Maximax:} choose \textbf{Large factory} ($A_1$), with best-case profit 132 thousand USD.\par
\textbf{Maximin:} choose \textbf{Compact line} ($A_3$), with worst-case profit 33 thousand USD.
\par\vspace{4mm}
\secInicio[sec:bicycle-assembly][sec:summarized-solutions]
\subsection{Solar Installation \hfill C4--2}
\label{summary:solar-installation}
Alternatives: $A_1$: Regional rollout; $A_2$: Town rollout; $A_3$: Village rollout; $A_4$: Service agreement.\par
States: $S_1$: Rapid adoption; $S_2$: Moderate adoption; $S_3$: Slow adoption.
\begin{center}
\begin{tabular}{@{}lcc@{}}\toprule
Alternative & Row maximum $M_i$ & Row minimum $m_i$\\\midrule
Regional rollout ($A_1$) & 168 & -39\\
Town rollout ($A_2$) & 108 & 17\\
Village rollout ($A_3$) & 73 & 31\\
Service agreement ($A_4$) & 35 & 35\\
\bottomrule\end{tabular}
\end{center}
\[\max_i M_i=\max\{168,108,73,35\}=168,\qquad \max_i m_i=\max\{-39,17,31,35\}=35.\]
\textbf{Maximax:} choose \textbf{Regional rollout} ($A_1$), with best-case profit 168 thousand USD.\par
\textbf{Maximin:} choose \textbf{Service agreement} ($A_4$), with worst-case profit 35 thousand USD.
\par\vspace{4mm}
\secInicio[sec:solar-installation][sec:summarized-solutions]
\subsection{Equipment Leasing \hfill C4--2}
\label{summary:equipment-leasing}
Alternatives: $A_1$: Buy large fleet; $A_2$: Buy medium fleet; $A_3$: Lease small fleet; $A_4$: Pool equipment.\par
States: $S_1$: Peak utilisation; $S_2$: Busy utilisation; $S_3$: Routine utilisation; $S_4$: Quiet utilisation.
\begin{center}
\begin{tabular}{@{}lcc@{}}\toprule
Alternative & Row maximum $M_i$ & Row minimum $m_i$\\\midrule
Buy large fleet ($A_1$) & 156 & -46\\
Buy medium fleet ($A_2$) & 113 & 11\\
Lease small fleet ($A_3$) & 77 & 25\\
Pool equipment ($A_4$) & 54 & 32\\
\bottomrule\end{tabular}
\end{center}
\[\max_i M_i=\max\{156,113,77,54\}=156,\qquad \max_i m_i=\max\{-46,11,25,32\}=32.\]
\textbf{Maximax:} choose \textbf{Buy large fleet} ($A_1$), with best-case profit 156 thousand USD.\par
\textbf{Maximin:} choose \textbf{Pool equipment} ($A_4$), with worst-case profit 32 thousand USD.
\par\vspace{4mm}
\secInicio[sec:equipment-leasing][sec:summarized-solutions]
\clearpage
\subsection{Festival Meals \hfill C5--2}
\label{summary:festival-meals}
Alternatives: $A_1$: Full menu; $A_2$: Balanced menu; $A_3$: Limited menu.\par
States: $S_1$: Crowded; $S_2$: Average attendance; $S_3$: Sparse attendance.
\[\boldsymbol{b}=\bigl(87,54,34\bigr),\qquad b_j=\max_k P(A_k,S_j).\]
All payoffs and regrets in this calculation are in thousand USD.
\begin{center}
\begin{tabular}{@{}lcccc@{}}\toprule
Alternative & $S_1$ & $S_2$ & $S_3$ & Maximum regret\\\midrule
$A_1$ & $\begin{gathered}87-87\\=0\end{gathered}$ & $\begin{gathered}54-37\\=17\end{gathered}$ & $\begin{gathered}34-(-16)\\=50\end{gathered}$ & 50\\
$A_2$ & $\begin{gathered}87-65\\=22\end{gathered}$ & $\begin{gathered}54-54\\=0\end{gathered}$ & $\begin{gathered}34-21\\=13\end{gathered}$ & 22\\
$A_3$ & $\begin{gathered}87-43\\=44\end{gathered}$ & $\begin{gathered}54-41\\=13\end{gathered}$ & $\begin{gathered}34-34\\=0\end{gathered}$ & 44\\
\bottomrule\end{tabular}
\end{center}
\[\min_i\max_j r(A_i,S_j)=\min\{50,22,44\}=22.\]
Choose \textbf{Balanced menu} ($A_2$): its maximum regret is \textbf{22 thousand USD}. 
\par\vspace{4mm}
\secInicio[sec:festival-meals][sec:summarized-solutions]
\subsection{Laundry Capacity \hfill C5--2}
\label{summary:laundry-capacity}
Alternatives: $A_1$: Build plant; $A_2$: Upgrade machines; $A_3$: Lease machines; $A_4$: Contract service.\par
States: $S_1$: Heavy workload; $S_2$: Routine workload; $S_3$: Light workload.
\[\boldsymbol{b}=\bigl(126,73,41\bigr),\qquad b_j=\max_k P(A_k,S_j).\]
All payoffs and regrets in this calculation are in thousand USD.
\begin{center}
\begin{tabular}{@{}lcccc@{}}\toprule
Alternative & $S_1$ & $S_2$ & $S_3$ & Maximum regret\\\midrule
$A_1$ & $\begin{gathered}126-126\\=0\end{gathered}$ & $\begin{gathered}73-47\\=26\end{gathered}$ & $\begin{gathered}41-(-32)\\=73\end{gathered}$ & 73\\
$A_2$ & $\begin{gathered}126-94\\=32\end{gathered}$ & $\begin{gathered}73-73\\=0\end{gathered}$ & $\begin{gathered}41-19\\=22\end{gathered}$ & 32\\
$A_3$ & $\begin{gathered}126-69\\=57\end{gathered}$ & $\begin{gathered}73-61\\=12\end{gathered}$ & $\begin{gathered}41-38\\=3\end{gathered}$ & 57\\
$A_4$ & $\begin{gathered}126-41\\=85\end{gathered}$ & $\begin{gathered}73-41\\=32\end{gathered}$ & $\begin{gathered}41-41\\=0\end{gathered}$ & 85\\
\bottomrule\end{tabular}
\end{center}
\[\min_i\max_j r(A_i,S_j)=\min\{73,32,57,85\}=32.\]
Choose \textbf{Upgrade machines} ($A_2$): its maximum regret is \textbf{32 thousand USD}. 
\par\vspace{4mm}
\secInicio[sec:laundry-capacity][sec:summarized-solutions]
\clearpage
\subsection{Juice Production \hfill C5--2}
\label{summary:juice-production}
Alternatives: $A_1$: Large batch; $A_2$: Staged batches; $A_3$: Small batch.\par
States: $S_1$: Exceptional sales; $S_2$: Brisk sales; $S_3$: Ordinary sales; $S_4$: Poor sales.
\[\boldsymbol{b}=\bigl(1080,730,510,390\bigr),\qquad b_j=\max_k P(A_k,S_j).\]
All payoffs and regrets in this calculation are in USD.
\begin{center}
\begin{tabular}{@{}lccccc@{}}\toprule
Alternative & $S_1$ & $S_2$ & $S_3$ & $S_4$ & Maximum regret\\\midrule
$A_1$ & $\begin{gathered}1080-1080\\=0\end{gathered}$ & $\begin{gathered}730-610\\=120\end{gathered}$ & $\begin{gathered}510-230\\=280\end{gathered}$ & $\begin{gathered}390-(-270)\\=660\end{gathered}$ & 660\\
$A_2$ & $\begin{gathered}1080-810\\=270\end{gathered}$ & $\begin{gathered}730-730\\=0\end{gathered}$ & $\begin{gathered}510-510\\=0\end{gathered}$ & $\begin{gathered}390-150\\=240\end{gathered}$ & 270\\
$A_3$ & $\begin{gathered}1080-530\\=550\end{gathered}$ & $\begin{gathered}730-550\\=180\end{gathered}$ & $\begin{gathered}510-480\\=30\end{gathered}$ & $\begin{gathered}390-390\\=0\end{gathered}$ & 550\\
\bottomrule\end{tabular}
\end{center}
\[\min_i\max_j r(A_i,S_j)=\min\{660,270,550\}=270.\]
Choose \textbf{Staged batches} ($A_2$): its maximum regret is \textbf{270 USD}. 
\par\vspace{4mm}
\secInicio[sec:juice-production][sec:summarized-solutions]
\clearpage
\subsection{Coffee Roasting \hfill C6--2}
\label{summary:coffee-roasting}
Alternatives: $A_1$: Specialty roast; $A_2$: Blended roast; $A_3$: Everyday roast.\par
States: $S_1$: Export boom; $S_2$: Regular exports; $S_3$: Export slump.
\[\boldsymbol{q}=(0.35,0.40,0.25),\qquad 0.35+0.40+0.25=1.\]
All probabilities lie in $[0,1]$, so they are nonnegative and valid.
\begin{align*}
EMV(A_1)&=0.35(98)+0.40(49)+0.25(-17)=49.65,\\
EMV(A_2)&=0.35(74)+0.40(63)+0.25(27)=57.85,\\
EMV(A_3)&=0.35(51)+0.40(53)+0.25(42)=49.55.
\end{align*}
\[\max\{49.65,57.85,49.55\}=57.85.\]
Choose \textbf{Blended roast} ($A_2$), with expected profit \textbf{57.85 thousand USD} (USD 57,850). 
Expected profit is not guaranteed profit.
\par\vspace{4mm}
\secInicio[sec:coffee-roasting][sec:summarized-solutions]
\subsection{Farm Visitor Facilities \hfill C6--2}
\label{summary:farm-visitor-facilities}
Alternatives: $A_1$: Visitor centre; $A_2$: Covered terrace; $A_3$: Picnic shelter; $A_4$: Lease to operator.\par
States: $S_1$: Busy season; $S_2$: Typical season; $S_3$: Quiet season.
\[\boldsymbol{q}=(0.25,0.45,0.30),\qquad 0.25+0.45+0.30=1.\]
All probabilities lie in $[0,1]$, so they are nonnegative and valid.
\begin{align*}
EMV(A_1)&=0.25(117)+0.45(43)+0.30(-29)=39.90,\\
EMV(A_2)&=0.25(89)+0.45(64)+0.30(23)=57.95,\\
EMV(A_3)&=0.25(63)+0.45(57)+0.30(40)=53.40,\\
EMV(A_4)&=0.25(32)+0.45(32)+0.30(32)=32.
\end{align*}
\[\max\{39.90,57.95,53.40,32\}=57.95.\]
Choose \textbf{Covered terrace} ($A_2$), with expected profit \textbf{57.95 thousand USD} (USD 57,950). 
Expected profit is not guaranteed profit.
\par\vspace{4mm}
\secInicio[sec:farm-visitor-facilities][sec:summarized-solutions]
\subsection{Translation Capacity \hfill C6--2}
\label{summary:translation-capacity}
Alternatives: $A_1$: Permanent staff; $A_2$: Staff plus freelancers; $A_3$: Core staff.\par
States: $S_1$: Order rush; $S_2$: Busy orders; $S_3$: Routine orders; $S_4$: Few orders.
\[\boldsymbol{q}=(0.15,0.35,0.25,0.25),\qquad 0.15+0.35+0.25+0.25=1.\]
All probabilities lie in $[0,1]$, so they are nonnegative and valid.
\begin{align*}
EMV(A_1)&=0.15(137)+0.35(81)+0.25(29)+0.25(-43)=45.40,\\
EMV(A_2)&=0.15(99)+0.35(86)+0.25(59)+0.25(17)=63.95,\\
EMV(A_3)&=0.15(65)+0.35(67)+0.25(61)+0.25(48)=60.45.
\end{align*}
\[\max\{45.40,63.95,60.45\}=63.95.\]
Choose \textbf{Staff plus freelancers} ($A_2$), with expected profit \textbf{63.95 thousand USD} (USD 63,950). 
Expected profit is not guaranteed profit.
\par\vspace{4mm}
\secInicio[sec:translation-capacity][sec:summarized-solutions]
\clearpage
\subsection{Mobile Game Release \hfill C7--3}
\label{summary:mobile-game-release}
Alternatives: $A_1$: Wide release; $A_2$: Limited release.\par
States: $S_1$: Viral uptake; $S_2$: Moderate uptake; $S_3$: Low uptake.
The mutually exclusive and exhaustive state probabilities sum to one:
\[P(S_3)=1-p-p/2=1-3p/2,\qquad\boldsymbol{q}=\left(p,\frac p2,1-\frac{3p}{2}\right).\]
Every probability must lie in $[0,1]$:
\[0\le p\le1,\quad 0\le p/2\le1\iff0\le p\le2,\quad 0\le1-3p/2\le1\iff0\le p\le\frac23.\]
Their intersection is $\boxed{0\le p\le \frac{2}{3}}$.
The remaining part of $[0,1]$, $\frac{2}{3}<p\le1$, is infeasible because $P(S_3)<0$. No decision is valid there.
Expected profits are in thousand USD:
\begin{align*}
E_1&=EMV(A_1)=108p+44(p/2)-10(1-3p/2)\\&=108p+22p-10+15p=-10+145p,\\
E_2&=EMV(A_2)=58p+64(p/2)+30(1-3p/2)\\&=58p+32p+30-45p=30+45p.
\end{align*}
Equating the two linear EMVs gives
\[E_1=E_2\iff 100p=40\iff\boxed{p=\frac{2}{5}}.\]
This is the only threshold, since the slopes differ, and $0<\frac{2}{5}<\frac{2}{3}$. At the tie, $E_1=E_2=48$ thousand USD.
The difference $E_1-E_2=-40+100p$ is negative below the threshold and positive above it.
At $p=0$, $(E_1,E_2)=(-10,30)$; at $p=\frac{2}{3}$, $(E_1,E_2)=(\frac{260}{3},60)$.
Both feasible endpoints agree with the rule:
\[\boxed{\begin{cases}
\text{Limited release }(A_2),&0\le p<\frac{2}{5},\\
\text{Both alternatives tied},&p=\frac{2}{5},\\
\text{Wide release }(A_1),&\frac{2}{5}<p\le\frac{2}{3}.
\end{cases}}\]
Either alternative is optimal at the tie. Expected profit is not guaranteed profit.
\par\vspace{4mm}
\secInicio[sec:mobile-game-release][sec:summarized-solutions]
\clearpage
\subsection{Timber Supply Agreement \hfill C7--3}
\label{summary:timber-supply-agreement}
Alternatives: $A_1$: Spot sales; $A_2$: Reserved supply.\par
States: $S_1$: Scarce timber; $S_2$: Balanced supply; $S_3$: Abundant timber.
The mutually exclusive and exhaustive state probabilities sum to one:
\[P(S_3)=1-2p-p=1-3p,\qquad\boldsymbol{q}=(2p,p,1-3p).\]
Every probability must lie in $[0,1]$:
\[0\le2p\le1\iff0\le p\le\frac12,\quad 0\le p\le1,\quad 0\le1-3p\le1\iff0\le p\le\frac13.\]
Their intersection is $\boxed{0\le p\le \frac{1}{3}}$.
The remaining part of $[0,1]$, $\frac{1}{3}<p\le1$, is infeasible because $P(S_3)<0$. No decision is valid there.
For $p>\frac12$, $P(S_1)>1$ as well.
Expected profits are in thousand USD:
\begin{align*}
E_1&=EMV(A_1)=30(2p)+66p+94(1-3p)\\&=60p+66p+94-282p=94-156p,\\
E_2&=EMV(A_2)=85(2p)+46p+40(1-3p)\\&=170p+46p+40-120p=40+96p.
\end{align*}
Equating the two linear EMVs gives
\[E_1=E_2\iff -252p=-54\iff\boxed{p=\frac{3}{14}}.\]
This is the only threshold, since the slopes differ, and $0<\frac{3}{14}<\frac{1}{3}$. At the tie, $E_1=E_2=\frac{424}{7}$ thousand USD.
The difference $E_1-E_2=54-252p$ is positive below the threshold and negative above it.
At $p=0$, $(E_1,E_2)=(94,40)$; at $p=\frac{1}{3}$, $(E_1,E_2)=(42,72)$.
Both feasible endpoints agree with the rule:
\[\boxed{\begin{cases}
\text{Spot sales }(A_1),&0\le p<\frac{3}{14},\\
\text{Both alternatives tied},&p=\frac{3}{14},\\
\text{Reserved supply }(A_2),&\frac{3}{14}<p\le\frac{1}{3}.
\end{cases}}\]
Either alternative is optimal at the tie. Expected profit is not guaranteed profit.
\par\vspace{4mm}
\secInicio[sec:timber-supply-agreement][sec:summarized-solutions]
}
\end{document}
